\section*{結言}
今回の観測では、フロンティアモデルの性能競争と、AIの利用をどう制御するかという課題が同時に前面へ出た。Gemini 4 Carbonはサイバー分野での能力を掲げ、OpenAIのModel Spec改訂は人間の意図や主体性を重視する方針を示した。いずれもモデルの出力性能だけでは測れない運用面の論点を含む。

開発・実装の側面では、音声転写のオープンウェイトモデル、端末内推論、エージェント向けターミナル、Copilot関連機能、APIの入力制限変更などが並んだ。モデルを単体で評価する段階から、既存ツールや業務フローに組み込むための基盤整備へ関心が広がっている。一方、各ツールの対応範囲や性能は発表側の説明に依存する部分が残る。

MetaのモデルIDやブラウザセキュリティ研究、オープンウェイト採用の議論には、公式確認や実証が不足している。観測窓より前のCarbon発表と、窓内に続いた反応も区別する必要がある。次に注目したいのは、新モデルの実際の提供範囲、エージェント基盤の運用実績、未確認の主張に対する一次資料の公開である。
